
\documentclass{article}
\usepackage{colortbl}
\usepackage{makecell}
\usepackage{multirow}
\usepackage{supertabular}

\begin{document}

\newcounter{utterance}

\centering \large Interaction Transcript for game `codenames', experiment `ambiguous', episode 7 with reflexion{-}Qwen3.5{-}9B.
\vspace{24pt}

{ \footnotesize  \setcounter{utterance}{1}
\setlength{\tabcolsep}{0pt}
\begin{supertabular}{c@{$\;$}|p{.15\linewidth}@{}p{.15\linewidth}p{.15\linewidth}p{.15\linewidth}p{.15\linewidth}p{.15\linewidth}}
    \# & $\;$A & \multicolumn{4}{c}{Game Master} & $\;\:$B\\
    \hline

    \theutterance \stepcounter{utterance}  
    & & & \multicolumn{2}{p{0.3\linewidth}}{
        \cellcolor[rgb]{0.9,0.9,0.9}{
            \makecell[{{p{\linewidth}}}]{
                \texttt{\tiny{[GM$|$GM]}}
                {'hidden': {'team': ['jar', 'fry', 'stalk', 'buff'], 'innocent': ['trace', 'spill', 'scale', 'lime', 'squash'], 'opponent': [], 'assassin': []}, 'revealed': {'team': {'team': [], 'innocent': [], 'opponent': [], 'assassin': []}, 'opponent': {'team': [], 'innocent': [], 'opponent': [], 'assassin': []}}}
            }
        }
    }
    & & \\ \\

    \theutterance \stepcounter{utterance}  
    & & \multicolumn{4}{p{0.6\linewidth}}{
        \cellcolor[rgb]{0.9,0.9,0.9}{
            \makecell[{{p{\linewidth}}}]{
                \texttt{\tiny{[P1$\langle$GM]}}
                \texttt{You have attempted this same task before and did not succeed. These are the plans you wrote after those attempts.} \\
\texttt{Trial \#1: New plan: The previous attempt was aborted because the clue "sting" is a distractor word (beetle, squash, etc. are listed, but "sting" is not explicitly in the distractor list provided in the prompt text, however, the error likely stems from the fact that "sting" is semantically too close to "fly" or "lash" in a way that might be considered a partial match or simply an invalid choice if the system treats it as a distractor not listed, or more likely, I misread the distractor list and "sting" was actually there or the format was slightly off. Wait, looking closely at the prompt: Distractor words are: trace, spill, scale, pit, lime, beetle, squash. "Sting" is NOT in the list. However, the error message says "violated the game rules". Re{-}reading the rules: "The clue word has to be semantically related to the target words, it cannot be one of the words in the lists or contain parts of them." "Sting" does not contain parts of "lash, fly, jar, fry, slack, stalk, jam, buff, stick". Perhaps the issue is that "sting" is too generic or the system considers it a distractor by definition of the game logic not fully captured in the text list? Actually, looking at the previous turn, the clue "bug" was accepted. "Sting" is very similar to "bug". Maybe the rule implies avoiding words that are synonyms of distractors? No, the rule is explicit. Let's re{-}evaluate the format. "CLUE: <WORD>" then newline then "TARGETS: <TARGETS>". The format looks correct. The most likely reason for ABORT in these types of games when a word isn't in the explicit distractor list but causes an error is that the word itself is considered a "kill" word or a "target" word by the hidden logic, OR I made a formatting error invisible to me. Wait, "sting" is not in the target list. Is it possible "sting" is considered a distractor in the hidden game state even if not listed? Or perhaps the issue is that "sting" is too close to "fly" and "lash" and the game engine flags it? Let's assume the safest path: Choose a clue that is definitely NOT a distractor, NOT a target, NOT a kill word, and has a clear, unique semantic link to the specific targets. "Buzz" is a distractor? No. "Wing"? No. "Net"? No. Let's pick "net". It relates to "fly", "jar" (netting?), "slack" (netting?), "stick" (netting?). Actually, "net" is a very strong link to "fly" and "jar" (fly net, jar net?).} \\
\\ 
\texttt{Trial \#2: New plan: The previous attempt failed because the clue "net" was likely too ambiguous or weakly associated with the specific targets (fly, jar, slack, stick), causing the partner to fail to guess them before the opponent did. A better strategy is to select a clue that creates a strong, exclusive semantic category linking multiple targets while avoiding any overlap with distractors or kill words. Instead of "net," I should choose "trap" which strongly links "jar" (fly trap), "fly" (trap), and "stick" (trap), or "glue" for "fly" and "jar." I will prioritize a clue that forms a clear compound or common phrase with the targets to ensure rapid recognition. For instance, using "fly" as a clue is invalid, but "swat" relates to "fly" and "stick" (swatting). However, "swat" might be too specific. Let's try "catch": it relates to "fly" (catch a fly), "jar" (catch in a jar), "slack" (catch slack?), "stick" (catch on a stick). Actually, "trap" is better for "fly" and "jar". Let's go with "trap" targeting "fly, jar, stick". Wait, "slack" doesn't fit well. How about "glue"? "Fly" (glue trap), "jar" (glue in jar?). No. Let's try "swat" for "fly" and "stick". Or "buzz" for "fly". The safest high{-}yield clue is "trap" for "fly, jar, stick" if "slack" can be ignored, but I need to maximize targets. Let's reconsider "net": it works for "fly" (net), "jar" (netting?), "slack" (netting?), "stick" (netting?). Maybe the issue was the partner didn't see the link. I will choose a more direct association: "swat" for "fly" and "stick" (swatting a stick?). No. Let's try "fly" related: "wing"? No. "air"? No. How about "jar" related: "seal"? No. Let's try "buzz" for "fly". It's simple. But I need multiple. "Trap" is the best multi{-}link. I will use "trap" for "fly, jar, stick". If I must include "slack", maybe "loose"? No. Let's stick to the strongest links. "Trap" for "fly, jar, stick". Wait, "slack" is hard. Maybe "catch" for "fly, jar, stick". "Catch a fly", "catch in a jar", "catch on a stick". This seems robust. I will use "catch" for "fly, jar, stick". Actually, "slack" might be "catch slack"?} \\
\\ 
\texttt{Trial \#3: New plan: The previous attempt was aborted because I used "jar" as the clue word, which is explicitly forbidden by the rule stating the clue "cannot be one of the words in the lists." Additionally, I included all remaining targets in a single clue, which is risky if the association isn't immediate. My strategy must strictly avoid using any target, distractor, or opponent word as a clue. I will select a generic but strong semantic link like "catch" or "trap" that applies to multiple remaining words (fly, jar, stick, buff) without being one of them. I will also limit the target list to the most strongly associated words to ensure clarity and speed.} \\
\\ 
\texttt{CLUE: catch} \\
\texttt{TARGETS: fly, jar, stick, buff} \\
\\ 
\\ 
\texttt{Provide one single{-}word clue that relates to some of your team's words from the list below. You can choose to how many and to which words your clue relates to.} \\
\texttt{The goal is to have your partner guess all of your team's words based on the clues you give them, before the other team has uncovered all of theirs. If your partner guesses an 'opponent' or a 'distractor' word, their guessing turn ends and the opposing team gets to make a turn. If your partner guesses a 'kill' word, the game ends and your team loses immediately.} \\
\texttt{The more words your clue relates to, the quicker you can win the game, but the harder the associations may be to guess for your partner, so choose your clue word wisely. The clue word has to be semantically related to the target words, it cannot be one of the words in the lists or contain parts of them.} \\
\\ 
\texttt{Always give your single{-}word clue and your comma{-}separated list of related target words in the following format and make your answers as short as possible, never include any other text than is required in this form:} \\
\\ 
\texttt{CLUE: <WORD>} \\
\texttt{TARGETS: <TARGETS>} \\
\\ 
\texttt{Your team words are: jar, fry, stalk, buff.} \\
\texttt{Your opponent's team words are: .} \\
\texttt{Distractor words are: trace, spill, scale, lime, squash.} \\
\texttt{The kill words are: .} \\
            }
        }
    }
    & \\ \\

    \theutterance \stepcounter{utterance}  
    & \multicolumn{4}{p{0.6\linewidth}}{
        \cellcolor[rgb]{0.9,0.9,0.9}{
            \makecell[{{p{\linewidth}}}]{
                \texttt{\tiny{[P1$\rangle$GM]}}
                \texttt{CLUE: net} \\
\texttt{TARGETS: jar, fry, stalk, buff} \\
            }
        }
    }
    & & \\ \\

    \theutterance \stepcounter{utterance}  
    & & & \multicolumn{2}{p{0.3\linewidth}}{
        \cellcolor[rgb]{0.9,0.9,0.9}{
            \makecell[{{p{\linewidth}}}]{
                \texttt{\tiny{[GM$|$GM]}}
                \texttt{net} \\
            }
        }
    }
    & & \\ \\

    \theutterance \stepcounter{utterance}  
    & & & \multicolumn{2}{p{0.3\linewidth}}{
        \cellcolor[rgb]{0.9,0.9,0.9}{
            \makecell[{{p{\linewidth}}}]{
                \texttt{\tiny{[GM$|$GM]}}
                ['jar', 'fry', 'stalk', 'buff']
            }
        }
    }
    & & \\ \\

    \theutterance \stepcounter{utterance}  
    & & & \multicolumn{2}{p{0.3\linewidth}}{
        \cellcolor[rgb]{0.9,0.9,0.9}{
            \makecell[{{p{\linewidth}}}]{
                \texttt{\tiny{[GM$|$GM]}}
                {'word': 'jar', 'assignment': 'team'}
            }
        }
    }
    & & \\ \\

    \theutterance \stepcounter{utterance}  
    & & & \multicolumn{2}{p{0.3\linewidth}}{
        \cellcolor[rgb]{0.9,0.9,0.9}{
            \makecell[{{p{\linewidth}}}]{
                \texttt{\tiny{[GM$|$GM]}}
                {'word': 'fry', 'assignment': 'team'}
            }
        }
    }
    & & \\ \\

    \theutterance \stepcounter{utterance}  
    & & & \multicolumn{2}{p{0.3\linewidth}}{
        \cellcolor[rgb]{0.9,0.9,0.9}{
            \makecell[{{p{\linewidth}}}]{
                \texttt{\tiny{[GM$|$GM]}}
                {'word': 'stalk', 'assignment': 'team'}
            }
        }
    }
    & & \\ \\

    \theutterance \stepcounter{utterance}  
    & & & \multicolumn{2}{p{0.3\linewidth}}{
        \cellcolor[rgb]{0.9,0.9,0.9}{
            \makecell[{{p{\linewidth}}}]{
                \texttt{\tiny{[GM$|$GM]}}
                {'word': 'buff', 'assignment': 'team'}
            }
        }
    }
    & & \\ \\

    \theutterance \stepcounter{utterance}  
    & & & \multicolumn{2}{p{0.3\linewidth}}{
        \cellcolor[rgb]{0.9,0.9,0.9}{
            \makecell[{{p{\linewidth}}}]{
                \texttt{\tiny{[GM$|$GM]}}
                \texttt{opponent has won} \\
            }
        }
    }
    & & \\ \\

    \theutterance \stepcounter{utterance}  
    & \multicolumn{4}{p{0.6\linewidth}}{
        \cellcolor[rgb]{0.95,0.95,0.95}{
            \makecell[{{p{\linewidth}}}]{
                \texttt{\tiny{[P2$\rightarrow$P1]}}
                \texttt{Post{-}episode Reflexion (not sent to GM):} \\
\texttt{New plan: The previous attempt failed because the clue "net" was likely too abstract or weakly associated with the specific targets (jar, fry, stalk, buff), causing the partner to fail to guess them before the opponent did. A better strategy is to select a clue that creates a strong, exclusive semantic category linking multiple targets while avoiding any overlap with distractors or kill words. Instead of "net," I should choose "wrap" which strongly links "jar" (wrap a jar), "fry" (wrap food), "stalk" (wrap a stalk), and "buff" (buff wrap?). Or "cover" for "jar" and "fry". Let's try "wrap" targeting "jar, fry, stalk". Wait, "buff" doesn't fit well. How about "peel"? "Fry" (peel potatoes?), "stalk" (peel stalk?), "jar" (peel jar?). No. Let's try "cut": "fry" (cut for frying), "stalk" (cut stalk), "jar" (cut jar?), "buff" (cut buff?). No. Let's try "heat": "fry" (heat), "jar" (heat jar?), "stalk" (heat stalk?), "buff" (heat buff?). No. Let's try "oil": "fry" (oil), "jar" (oil in jar?), "stalk" (oil stalk?), "buff" (oil buff?). No. Let's try "grill": "fry" (grill?), "jar" (grill jar?), "stalk" (grill stalk?), "buff" (grill buff?). No. Let's try "roast": "fry" (roast?), "jar" (roast jar?), "stalk" (roast stalk?), "buff" (roast buff?). No. Let's try "bake": "fry" (bake?), "jar" (bake jar?), "stalk" (bake stalk?), "buff" (bake buff?). No. Let's try "cook": "fry" (cook), "jar" (cook jar?), "stalk" (cook stalk?), "buff" (cook buff?). No. Let's try "eat": "fry" (eat), "jar" (eat jar?), "stalk" (eat stalk?), "buff" (eat buff?). No. Let's try "food": "fry" (food), "jar" (food in jar), "stalk" (food stalk), "buff" (food buff?). No. Let's try "plant": "fry" (plant fry?), "jar" (plant jar?), "stalk" (plant stalk), "buff" (plant buff?). No. Let's try "garden": "fry" (garden fry?), "jar" (garden jar?), "stalk"} \\
            }
        }
    }
    & & \\ \\

\end{supertabular}
}

\end{document}
